% ============================================================================
% Section 5: Fault Diagnosis (Task 5)
% Handout s10.2 requires ONLY: expected-versus-observed diagnosis table,
% faulty waveform, corrected waveform, register field responsible, short
% explanation of the correction. Both captures should be the SAME measurement
% (RUN_TASK 5 vs RUN_TASK 4).
%
% Still to do: figures/task5_faulty.jpeg and figures/task5_corrected.jpeg
% ============================================================================
\section{Fault Diagnosis}
\label{sec:fault}

\begin{table}[htbp]
  \centering
  \caption{Fault diagnosis (\reg{RUN\_TASK}~5 against \reg{RUN\_TASK}~4).}
  \label{tab:fault}
  \small
  \begin{tabularx}{\linewidth}{@{}XXXX@{}}
    \toprule
    \textbf{Expected} & \textbf{Observed} & \textbf{Cause} & \textbf{Correction} \\
    \midrule
    Mode~(0,0): MOSI changes on falling SCK edges and is stable at each rising
    edge; status 0x00, read 0x44, write done in \SI{3}{\milli\second}
    & MOSI changes at the rising SCK edges; status and read-back 0xFF; write
    poll times out after \SI{50}{\milli\second}
    & \reg{SPI2\_CR1} = 0x0365: CPHA (bit 0) = 1, i.e.\ mode~(0,1); MOSI changes on rising edges
    & CPHA = 0, \reg{SPI2\_CR1} = 0x0364: mode~(0,0) (\reg{RUN\_TASK}~4) \\
    \bottomrule
  \end{tabularx}
\end{table}

\begin{figure}[htbp]
  \centering
  \begin{subfigure}[t]{0.48\linewidth}
    \centering
    \subscope{task5_faulty.jpeg}{Faulty (\reg{RUN\_TASK}~$=5$).}{fig:fault-bad}
  \end{subfigure}\hfill
  \begin{subfigure}[t]{0.48\linewidth}
    \centering
    \subscope{task5_corrected.jpeg}{Corrected (\reg{RUN\_TASK}~$=4$).}{fig:fault-good}
  \end{subfigure}
  \caption{The same read-only transaction, faulty and corrected:
    CH1 = SCK, CH2 = MOSI.}
  \label{fig:fault}
\end{figure}

\paragraph{Register field responsible.}
\reg{SPI2\_CR1} bit~0, CPHA: 1 $\rightarrow$ 0 (CR1 0x0365 $\rightarrow$ 0x0364; CR2 = 0x1700 in both builds; CPOL = 0 in both, so SCK idles low in both)~\cite[\S28.9.1, \S28.5.6]{rm0091}.

\paragraph{Correction.}
The EEPROM latches SI on the rising edge of SCK and needs the data stable
around that edge~\cite[p.~5, Tab.~7]{eeprom}. With CPHA~=~1 the STM32 samples
on the second (falling) edge and therefore shifts MOSI on the first (rising)
edge~\cite[\S28.5.6]{rm0091}, so the EEPROM latched incorrect opcode bits, ignored every
command and left SO tri-stated; the MISO pull-up then returned 0xFF, and RDY
appeared permanently set. Restoring CPHA to 0 returns
the interface to mode~(0,0), so MOSI is stable at every rising edge
(\cref{fig:fault-good}): the EEPROM decodes the commands, the status reads
0x00 and 0x44 is read back from $A$.
